
\begin{exercise}[Discontinuity points of distribution functions \Coffeecup\Coffeecup]
    Let \(X\) be a random variable and \(F(x) = \prob{X \le x}\) its
    distribution function. Show that \(F\) is continuous at \(x\) if and only if
    \(\prob{X = x} = 0\).
    \textbf{Hint:} Use Exercise \ref{ex: monotone functions have countably many discontinuities}.
\end{exercise}
\begin{solution}
    Since the sets \(\set{X \le x - \frac1n}\) are increasing
    in size, we have by continuity of the probability measure
    \[
        \prob{X < x}
        = \prob[\Big]{\bigcup_{n\in \nat} \set{X \le x - \tfrac1n}}
        = \lim_{n\to\infty} \prob{X \le x - \tfrac1n}
        = F_-(x).
    \]
    Consequently,
    \[
        \prob{X = x} = \prob{X \le x} - \prob{X < x} = F(x) - F_-(x).
    \]
    By Exercise \ref{ex: monotone functions have countably many
    discontinuities} \ref{it: F continuous if and only if all coincide} we
    have that \(F\) is continuous at \(x\) if and only if \(F(x) = F_-(x)\).
    Consequently continuity of \(F\) is equivalent to \(\prob{X = x} = 0\).
\end{solution}
